\section{Small-window information along realizable channels}
\label{sec:v53-window-information}

The window inverse and the physical sampling problem use the same boundary
measure, but they do not preserve the same information.  We make this
distinction quantitative within a family of actual periodic tables.  The
small-window calculation for quadratic functional actions was communicated
in~\cite{A2ReviewV52}.  The argument below needs only the leading Taylor
coefficient of an actual action.  It realizes the effect without requiring
an exactly quadratic billiard action, gives a matching two-point testing
scale, and transfers that experiment to sufficiently long finite bridges.
The recording profiles are fixed as the window shrinks.  They may differ
between the two alternatives, as permitted by the nuisance model of
Theorem~\ref{thm:v52-window-extraction}.

\subsection{A normalized window calculation}

Fix $d>0$.  Let $S_i$, $i\in\{0,1\}$, be even $C^4$ functions near
zero, with $S_i(0)=S_i'(0)=0$, $S_i''>0$, and put $a_i=S_i''(0)$.
For a sufficiently small $h>0$ take the window $(-h,h)^2$ and the
separate effective factors $(d-S_i(u))^{-1}$ and $(d-S_i(v))^{-1}$.
Positive constants multiplying either factor do not change the conditional
law.  After the known rescaling $u=hz$, $v=hw$, write $Q_{i,h}$ for
that law on $[-1,1]^2$, ignoring edges of measure zero.  Define
\begin{equation}\label{eq:v53-q}
 t_i(u)=\frac{S_i(u)}{d-S_i(u)},\qquad
 I_i(h)=\int_{-1}^1t_i(hz)\,\dd z,\qquad
 q_{i,h}(z,w)=\frac{1-t_i(hz)t_i(hw)}{4-I_i(h)^2}.
\end{equation}
Indeed the unnormalized window density is
$d^{-1}(1-t_i(u)t_i(v))$.  Its integral includes the term $I_i(h)^2$;
discarding this term changes the leading information coefficient.
Throughout this section the squared Hellinger distance is
\[
 H^2(P,Q)=\int(\sqrt p-\sqrt q)^2,
 \qquad \operatorname{TV}(P,Q)=\tfrac12\int|p-q|.
\]

\begin{lemma}[Window expansion and quadratic information loss]
\label{lem:v53-window-expansion}
Put $F(z,w)=1/9-z^2w^2$.  Uniformly on the rescaled square,
\begin{equation}\label{eq:v53-density-expansion}
 q_{i,h}=\frac14+\frac{a_i^2h^4}{16d^2}F+O(h^6).
\end{equation}
If $a_0\ne a_1$, then
\begin{align}
 \operatorname{TV}(Q_{0,h},Q_{1,h})&=\Theta(h^4),
                                             \label{eq:v53-tv}\\
 H^2(Q_{0,h},Q_{1,h})
  &=\frac{7(a_1^2-a_0^2)^2}{16200d^4}h^8+O(h^{10}),
                                             \label{eq:v53-hellinger}\\
 \mathbb E_{Q_{i,h}}F
  &=\frac{14a_i^2}{2025d^2}h^4+O(h^6),\qquad
 \operatorname{Var}_{Q_{i,h}}F=\frac{56}{2025}+O(h^4).
                                             \label{eq:v53-moments}
\end{align}
All remainder bounds are uniform on families with a common $C^4$ bound,
positive $a_i$, and a common positive window margin.
\end{lemma}

\begin{proof}
Evenness and Taylor's theorem give
$S_i(hz)=a_ih^2z^2/2+O(h^4)$, uniformly for $|z|\le1$.
Expansion of the reciprocal on a fixed positive neighborhood gives
$t_i(hz)=a_ih^2z^2/(2d)+O(h^4)$ and
$I_i(h)=a_ih^2/(3d)+O(h^4)$.  Expanding the numerator and denominator
of \eqref{eq:v53-q} proves \eqref{eq:v53-density-expansion}.
The square densities are bounded away from zero for small $h$.
Consequently
\[
 \sqrt{q_{1,h}}-\sqrt{q_{0,h}}
 =\frac{q_{1,h}-q_{0,h}}{\sqrt{q_{1,h}}+\sqrt{q_{0,h}}}
 =\frac{(a_1^2-a_0^2)h^4}{16d^2}F+O(h^6).
\]
Here the limiting denominator is one.  Direct integration gives
\begin{equation}\label{eq:v53-F-integrals}
 \int_{[-1,1]^2}F=0,\qquad
 \int_{[-1,1]^2}F^2
 =4\left(\frac1{25}-\frac1{81}\right)=\frac{224}{2025}.
\end{equation}
Squaring and integrating proves \eqref{eq:v53-hellinger}.
The nonzero continuous function $F$ also has a positive integral of its
absolute value; this proves \eqref{eq:v53-tv}.  Integrating $Fq_{i,h}$
and $F^2q_{i,h}$ proves \eqref{eq:v53-moments}.  The stated uniformity
follows from the common Taylor and reciprocal bounds.  No higher action
coefficient is set equal to zero in this proof.
\end{proof}

\subsection{Actual tables and the attainable discrimination scale}

The model below fixes the lattice and all channel gaps.  The only changing
geometric quantity is the obstacle shape.  Thus a gap observation or a
known origin cannot eliminate the information loss by itself.

\begin{theorem}[A realizable small-window experiment]
\label{thm:v53-realized-experiment}
There is a common-strip analytic-support family of noncircular periodic
tables with two fixed marked clear channels, the same gaps, and fixed
signed contact origins, with the following properties.
For each table there are positive smooth separate endpoint recording
probabilities, independent of $h$, such that all four same-type limiting
window laws are $Q_{s,h}$ of \eqref{eq:v53-q}, with the actual channel
action $S_s$.  Every fixed $h>0$ sufficiently small retains the exact
geometric fiber of Corollary~\ref{cor:v52-window-geometry}.

More explicitly, let $R=1/10$, let $s\in[R/128,R/64]$, put $L=I_2$, and
let the single obstacle orbit have support function
\begin{equation}\label{eq:v53-table-family}
 k_s(\theta)=R+s(\cos4\theta-1).
\end{equation}
The channels have deck marks $e_1,e_2$ and gap $g=4/5$.  The two types
and the two channels have the same even action, with
\begin{equation}\label{eq:v53-actual-quadratic}
 \kappa_s=(R-16s)^{-1},\qquad
 a_s=S_s''(0)=\sqrt{\kappa_s^2+2\kappa_s/g}.
\end{equation}
For the distinct alternatives $s_0=R/128$, $s_1=R/64$ this gives
\[
 a_0^2=7800/49,\qquad a_1^2=1900/9,\qquad
 \Delta=a_1^2-a_0^2=22900/441>0.
\]
Fix one common sufficiently small positive $d$.  Let $n=n_h$ independent
successful window records be allocated in any fixed way among the four
same-type laws.  The gaps and channel/type labels may also be retained.
Define
\begin{equation}\label{eq:v53-J}
             \mathcal J=\frac{7\Delta^2}{4050d^4}>0.
\end{equation}
For the two specified table/profile alternatives the following hold:
\begin{enumerate}
\item If $nh^8\to0$, their product total-variation distance tends to zero,
and the minimum equal-prior error tends to $1/2$.
\item If $nh^8\to\lambda\in(0,\infty)$, the minimum equal-prior error
converges to $\Phi(-\sqrt{\lambda\mathcal J}/2)$, where $\Phi$ is the
standard normal distribution function.
\item If $nh^8\to\infty$, a test of the empirical mean of $F$ has error
at most $\exp(-cnh^8)$ for some $c>0$ and all sufficiently small $h$.
The same mean test attains the limit in the second assertion when its
threshold is the midpoint of its two exact means.
\end{enumerate}
Thus $h^{-8}$ is the attainable successful-record scale for this pair.
The upper assertion is a two-point assertion with specified alternatives,
not a uniform estimator over unrestricted recording profiles.
\end{theorem}

\begin{proof}
\emph{Geometry and actual actions.}
The function in \eqref{eq:v53-table-family} is entire and real on the
real axis, and
\[
 k_s+k_s''=R-s-15s\cos4\theta\ge R-16s\ge3R/4>0.
\]
It is the support of a strictly convex noncircular body $K_s$, centered
at zero and contained in the disk of radius $R$, since $k_s\le R$.
Distinct $s$ give noncongruent bodies: their perimeters are $2\pi(R-s)$,
by integrating the radius of curvature.  Nonzero integer center
displacements have length at least one, so lifted closures are disjoint.
At the horizontal and vertical normals, $k_s=R$ and $k_s'=0$.
The supporting endpoints therefore give the unique closest pairs in
$K_s,K_s+e_i$, with gap $1-2R=4/5$.  Every other integer center is at
distance at least one from the associated unit axis segment.  Every
third body is at distance at least $1-R$ from that segment, and the
closest-contact segment is a subset.  Common small collars are clear.
No finite-horizon hypothesis is needed for these selected channels.

The quarter-turn symmetry makes the two channels congruent.  Reflection
in a channel axis sends every transverse coordinate to its negative;
the local flight action is invariant under this simultaneous sign change.
Uniqueness of the small half-line stationary sequence in
Lemma~\ref{lem:v4-halfline} therefore gives $S_s(-u)=S_s(u)$.
The reflection exchanging the two identical obstacles similarly identifies
the two types and their amplitudes.  These reflection arguments establish
symmetry of the functions of a fixed channel; they do not quotient an
unknown table by an improper Euclidean motion.  The family has a uniform
analytic collar and fixed finite smooth bounds by the same half-line
lemma.  With equal contact curvatures, $c_0=c_1=1+g\kappa_s$ and
$\cosh\gamma_s=1+g\kappa_s$; substituting in
\eqref{eq:v4-boundary-hessian} gives \eqref{eq:v53-actual-quadratic}.
This obtains the Taylor coefficient from an actual billiard, without
assuming its higher coefficients vanish.

\emph{Recording.}
Choose $d$ in the common positive-offset range and then $h_0>0$ inside
the common collar with $\max_{|u|\le2h_0}S_s(u)<d/4$ uniformly in $s$.
Let $B_s$ be the positive actual half-line amplitude.  On $[-h_0,h_0]$
choose both recording probabilities to be
\begin{equation}\label{eq:v53-efficiency}
              e_s(u)=\frac{c_*}{B_s(u)(d-S_s(u))}.
\end{equation}
A sufficiently small $c_*>0$ works simultaneously over the compact family.
For an explicit extension, take a smooth cutoff $\chi$ equal to one on
$[-h_0,h_0]$ and supported in $(-2h_0,2h_0)$, and set
$e_s=c_*\exp[-\chi\log(B_s(d-S_s))]$ on the gate interval, with the
cutoff product extended by zero.  The bounded exponent and a smaller
$c_*$ ensure $0<e_s\le1$ everywhere.  These profiles do not depend on $h$.
Independent recording of the two endpoints cancels $B_s(u)B_s(v)$ and
leaves exactly $c_*^2d^{-1}(1-t_s(u)t_s(v))$ on $(-h,h)^2$.
Conditioning and rescaling yield \eqref{eq:v53-q}.  The fixed origins
are zero, the cap is strictly positive on the window, and its critical
lines remain inside.  The window theorem therefore applies for each
$h>0$.  The profiles are admissible nuisance alternatives, not profiles
estimated or programmed from an unknown table by this theorem.

\emph{Testing.}
Symmetry makes every label's successful law the same $Q_{i,h}$ in these
coordinates; fixed label allocations add no distinction between the
alternatives.  Write $p_h=q_{0,h}$, $q_h=q_{1,h}$ and
$r_h=(q_h-p_h)/p_h$.  Lemma~\ref{lem:v53-window-expansion} gives
\begin{equation}\label{eq:v53-likelihood-expansion}
 \|r_h\|_\infty=O(h^4),\qquad
 \mathbb E_{p_h}r_h=0,\qquad
 \mathbb E_{p_h}r_h^2=\mathcal J h^8+O(h^{10}).
\end{equation}
For $nh^8\to0$, affinity tensorization gives
$H^2(p_h^{\otimes n},q_h^{\otimes n})\le nH^2(p_h,q_h)$;
Cauchy--Schwarz gives $\operatorname{TV}\le H$ in our convention.
The equal-prior optimal error is $(1-\operatorname{TV})/2$, proving the
first assertion.

Suppose $nh^8\to\lambda\in(0,\infty)$ and write the log likelihood ratio
as $\Lambda_h=\sum_{j=1}^n\log(1+r_h(X_j))$.  Uniform Taylor expansion
and \eqref{eq:v53-likelihood-expansion} give
\[
 \mathbb E_{p_h}\log(1+r_h)=-\tfrac12\mathcal J h^8+O(h^{10}),\qquad
 \operatorname{Var}_{p_h}\log(1+r_h)=\mathcal J h^8+O(h^{10}).
\]
Indeed the cubic error is bounded by
$C\|r_h\|_\infty\mathbb E_{p_h}r_h^2=O(h^{12})$.
The centered summands are uniformly $O(h^4)$, so the Lindeberg condition
holds and
\[
 \Lambda_h\ \Longrightarrow\ N(-\lambda\mathcal J/2,\lambda\mathcal J)
                         \quad\hbox{under }p_h^{\otimes n}.
\]
Under $q_h$, integration against $(1+r_h)p_h$ reverses the leading mean
and leaves the leading variance unchanged; it gives
$N(\lambda\mathcal J/2,\lambda\mathcal J)$.  The likelihood-ratio test
at zero is optimal.  The limiting distributions are continuous there,
so its two error probabilities converge to
$\Phi(-\sqrt{\lambda\mathcal J}/2)$.

Finally set $\mu_i(h)=\mathbb E_{Q_{i,h}}F$.  By
\eqref{eq:v53-moments},
\[
 \mu_1(h)-\mu_0(h)=\frac{14\Delta}{2025d^2}h^4+O(h^6)>0.
\]
The statistic $F$ lies in $[-8/9,1/9]$.  Testing its mean at
$(\mu_0+\mu_1)/2$ has each error at most
$\exp[-n(\mu_1-\mu_0)^2/2]$, by the bounded-variable exponential
inequality.  For completeness, that inequality follows by bounding the
centered log moment generating function of a variable in an interval of
length one by $t^2/8$, multiplying the bounds for independent variables,
and minimizing in $t$.  This proves the third assertion.  In the critical
regime the bounded triangular-array central limit theorem and
\eqref{eq:v53-moments} give
\[
 \frac{\sqrt n(\mu_1-\mu_0)}{2\sqrt{56/2025}}
       \longrightarrow\frac{\sqrt{\lambda\mathcal J}}2,
\]
so this mean test also attains the optimal critical limit.  Its threshold
uses the exact means of the specified alternatives, not a leading-order
threshold claimed uniformly valid for arbitrarily large $n$.
\end{proof}

\subsection{Finite flights and the separately retained mass}

The preceding scale counts successful window records.  It neither counts
all preparations nor describes a transcript retaining failures.  The
relative forward estimate permits a precise finite-flight comparison while
keeping this distinction intact.

\begin{corollary}[Uniform small-window finite-flight transfer]
\label{cor:v53-finite-flight}
Choose the gates and recording profiles as in
Theorem~\ref{thm:v53-realized-experiment}.  Let $Q_{s,J,h}$ be the rescaled
recorded law for an even $J$-flight bridge at time $Jg+d$, conditional on
recording both endpoints in the window.  There are $J_0$, $C<\infty$ and
$\tau\in(0,1)$, independent of $0<h\le h_0$ and $s$ in the displayed
compact interval, such that
\begin{equation}\label{eq:v53-finite-flight}
       \|q_{s,J,h}-q_{s,h}\|_\infty\le C\tau^J,
                  \qquad J\ge J_0.
\end{equation}
After decreasing $h_0$ and increasing $J_0$, both densities are bounded
below by a common $m>0$ on the rescaled square.  In particular,
\begin{align}
 H^2(Q_{s,J,h},Q_{s,h})&\le C\tau^{2J},
                                    \label{eq:v54-one-record-H}\\
 \operatorname{TV}\!\left(\bigotimes_{k=1}^n Q_{s,J,h}^{(k)},
                           \bigotimes_{k=1}^n Q_{s,h}^{(k)}\right)
   &\le \min\{1,C\sqrt n\,\tau^J\}.
                                    \label{eq:v54-product-TV}
\end{align}
Here the superscripts allow any fixed allocation among the four labelled
same-type laws.  Thus all optimal-error limits in
Theorem~\ref{thm:v53-realized-experiment} hold for the finite-flight pair
if $n_h\tau^{2J_h}\to0$.  The finite-flight mean-test error is at most
\begin{equation}\label{eq:v54-error-transfer}
          \exp(-cn_hh^8)+C\sqrt{n_h}\,\tau^{J_h}.
\end{equation}
For that mean test alone there is also a direct conclusion: whenever
$\tau^{J_h}=o(h^4)$, its error is at most $\exp(-c'n_hh^8)$ for all
sufficiently small $h$, without a product-approximation hypothesis on
$n_h$.

Let $\pi_{s,J}$ be the original gated-bridge success probability, before
endpoint recording and cropping, and let $r_{s,J}(h)$ be the conditional
probability of the extra recording and crop given that bridge.  Then
\begin{equation}\label{eq:v53-recording-mass}
 c_-h^2\le r_{s,J}(h)\le c_+h^2,\qquad
 r_{s,\infty}(h)=\frac{c_*^2h^2}{dZ_s}\{4-I_s(h)^2\},
\end{equation}
where $Z_s$ is the normalizer of the full limiting same-type endpoint law.
For independent repeated preparations at a fixed design, the expected
number needed for $n$ recorded successes is exactly
$n/(\pi_{s,J}r_{s,J}(h))$.  This expectation is not a deterministic
preparation-cap guarantee.
\end{corollary}

\begin{proof}
Let $f_{s,J}$ and $f_s$ be the successful full endpoint densities in the
true signed contact chart, before recording.  The relative factorization
in Theorem~\ref{thm:v4-factorization}, with its positive fixed-offset
normalizer, gives
\[
 \|f_{s,J}-f_s\|_{C^0([-h_0,h_0]^2)}\le C\tau^J.
\]
This is the density comparison in the proof of
Corollary~\ref{cor:v26-finite-flight-inverse}; the compact analytic family
supplies its uniform constants.  Shrinking $h_0$ and increasing $J_0$
make both densities uniformly positive and bounded on this square.
The recording profiles and their reciprocals are bounded there as well.
Writing $w_s(u,v)=e_s(u)e_s(v)$ and
$A_{s,J}(h)=\int_{(-h,h)^2}w_sf_{s,J}$, we have
\[
 A_{s,J}(h)\asymp h^2,\qquad
 |A_{s,J}(h)-A_{s,\infty}(h)|\le Ch^2\tau^J.
\]
Subtract the rescaled quotients
\[
 q_{s,J,h}(z,w)=
 \frac{h^2w_s(hz,hw)f_{s,J}(hz,hw)}{A_{s,J}(h)}.
\]
The numerator error and normalizer have the same area factor.  Their
quotient proves \eqref{eq:v53-finite-flight}, with no inverse-area loss.
The same bounds give the common positive lower density bound.

For any such two densities $p,q$ on $[-1,1]^2$,
\[
 H^2(P,Q)=\int\frac{(p-q)^2}{(\sqrt p+\sqrt q)^2}
       \le\frac{1}{4m}\int(p-q)^2
       \le m^{-1}\|p-q\|_\infty^2.
\]
This proves \eqref{eq:v54-one-record-H}.  If
$\mathfrak a(P,Q)=\int\sqrt{pq}=1-H^2(P,Q)/2$, independence gives
$\mathfrak a(\bigotimes P_k,\bigotimes Q_k)=\prod_k\mathfrak a(P_k,Q_k)$.
The inequality $1-\prod_k(1-x_k)\le\sum_kx_k$, $0\le x_k\le1$,
therefore gives
$H^2(\bigotimes P_k,\bigotimes Q_k)\le\sum_k H^2(P_k,Q_k)$.
Together with $\operatorname{TV}\le H$, this proves
\eqref{eq:v54-product-TV}, uniformly in the fixed label allocation.
Expectations of tests taking values in $[0,1]$ differ by at most total
variation under each alternative.  Taking the infimum over tests proves
the assertion about optimal errors; applying the bound to the mean test
proves \eqref{eq:v54-error-transfer}.  The former sufficient condition
$n_h\tau^{J_h}\to0$ and bound $\exp(-cn_hh^8)+Cn_h\tau^{J_h}$ follow
as weaker consequences, since $n_h\ge1$ and $\tau^{J_h}\le1$.

For the direct mean-test bound, put
$D_h=\mu_1(h)-\mu_0(h)=c_0h^4+O(h^6)$, where
$c_0=14\Delta/(2025d^2)>0$ and the means are those of the two ideal
laws.  The boundedness of $F$ and \eqref{eq:v53-finite-flight} show that
every finite-flight mean differs from its corresponding ideal mean by
at most $C\tau^J$.  If $\tau^J=o(h^4)$, the fixed ideal midpoint is
at distance at least $D_h/4$ from each finite-flight mean, on the correct
side.  The bounded-variable inequality used in the preceding theorem,
applied directly to the independent finite-flight variables of range
length one, bounds each error by $\exp(-n_hD_h^2/8)$.  This is
$\exp(-c'n_hh^8)$ for small $h$.  It compares the performance of this
statistic, not the total variation of two growing product experiments.

Finally $r_{s,J}(h)=A_{s,J}(h)$.  Substitution of the limiting density
and \eqref{eq:v53-efficiency} gives the exact mass formula in
\eqref{eq:v53-recording-mass}; the positivity bounds give its two-sided
estimate.  Each independent preparation is accepted with probability
$\pi_{s,J}r_{s,J}(h)>0$.  The sum of $n$ independent geometric waiting
times has the displayed expectation.  No conditioning on completion
of a finite preparation cap is used in this calculation.
\end{proof}

The Hellinger sharpening and the distinction between product transfer and
direct testing were communicated in~\cite{A2ReviewV53}.  The choices are
joint: for example, taking the smallest admissible even $J$ with
$\tau^J\le h/\sqrt{n_h+1}$ ensures $n_h\tau^{2J}\to0$.  At the critical
scale $n_hh^8\asymp1$, it suffices that $\tau^{J_h}=o(h^4)$.
For any fixed $\varepsilon>0$, a sufficient explicit choice there is
\[
 J_h=2\left\lceil
       \frac{(4+\varepsilon)\log(1/h)}{2|\log\tau|}
      \right\rceil,
\]
with $J_h$ increased to $J_0$ if necessary.  This improves a sufficient
flight-length coefficient; no optimal geometric convergence exponent is
asserted.  In the supercritical regime a direct bound for the mean test
need not approximate its entire growing sample law.

For every fixed window these records still determine their signed action
germs and exact table fibers.  The attained two-point scale does not
supply a uniform noisy inverse over arbitrary profiles, a fixed finite
vector recovering all analytic tables, or a calibration rule for unknown
windows.  It is a quantitative consequence of the same relative boundary
law on an explicitly realized geometric family.  The association algebra
is the inherited one~\cite{Osius2009}; neither its cancellation nor the
normal approximation is offered as a new general association principle.

\subsection{Waiting records, coarsening and a deterministic cap}

We next retain the accepted mass of the same experiment, rather than
normalizing it away.  Fix one channel and one contact type in
Theorem~\ref{thm:v53-realized-experiment}, and retain its two specified
table/profile alternatives $s_0,s_1$.  Only the acceptance indicator is
recorded on a failed preparation; no raw failed-shot position or residual
time is available.  Independent preparations are made at one even flight
number $J_h\to\infty$, fixed before acquisition.  Put
\begin{equation}\label{eq:v54-accepted-probability}
 p_{i,h}=\pi_{s_i,J_h}r_{s_i,J_h}(h),\qquad
 \gamma_i=\operatorname{arcosh}(1+g\kappa_{s_i}),\qquad
 \rho_h=p_{1,h}/p_{0,h}.
\end{equation}
Let $W_h$ count preparations up to and including the first accepted
window record, and let $Y_h$ be that rescaled mark.  Denote the binary
experiments observing $Y_h$ and $(W_h,Y_h)$ by $\mathcal E_h$ and
$\mathcal F_h$, respectively.  For binary experiments
$\mathcal U=(U_0,U_1)$, $\mathcal V=(V_0,V_1)$, use
\[
 \delta(\mathcal U,\mathcal V)
   =\inf_K\max_{i=0,1}\operatorname{TV}(KU_i,V_i),
\]
where $K$ ranges over Markov kernels independent of the unknown
alternative.  The repeat-to-first-success coarsening calculation below
was communicated in~\cite{A2ReviewV53}.  The deterministic-cap conclusion
also identifies the testing scale when the experiment can end without an
accepted mark.

\begin{theorem}[Accepted marks and charged first-acceptance records]
\label{thm:v54-waiting-and-cap}
For the specified pair and any even $J_h\to\infty$, as $h\downarrow0$,
\begin{equation}\label{eq:v54-escape-separation}
 p_{i,h}\asymp h^2e^{-\gamma_iJ_h},\qquad
 \gamma_1>\gamma_0>0,\qquad
 \rho_h\asymp e^{-(\gamma_1-\gamma_0)J_h}\longrightarrow0.
\end{equation}
No relation between $h$ and $\tau^{J_h}$ is needed for the following
conclusions.
\begin{enumerate}
\item The optimal equal-prior error for one accepted mark tends to
$1/2$, whereas that for the record $(W_h,Y_h)$ tends to zero.  More
precisely,
\begin{equation}\label{eq:v54-deficiency}
 \delta(\mathcal F_h,\mathcal E_h)=0,\qquad
 \delta(\mathcal E_h,\mathcal F_h)\longrightarrow\tfrac12.
\end{equation}
\item Fix an integer cap $b_h\ge1$.  Observe $(W_h,Y_h)$ if $W_h\le b_h$
and a single cemetery outcome $\partial$ otherwise.  Let $R_h^*(b_h)$
be the optimal equal-prior error in this capped experiment.  If
$b_hp_{0,h}\to\lambda\in[0,\infty)$, then
\begin{equation}\label{eq:v54-capped-critical-risk}
                       R_h^*(b_h)\longrightarrow\tfrac12e^{-\lambda}.
\end{equation}
The test selecting alternative zero exactly when an acceptance occurs
before the cap attains this limit.  For arbitrary cap sequences,
$R_h^*(b_h)\to0$ if and only if $b_hp_{0,h}\to\infty$.
\item The stopped experiment uses at most $b_h$ preparations.  Its
expected charge under alternative $i$ is exactly
\begin{equation}\label{eq:v54-capped-charge}
 \mathbb E_i\min(W_h,b_h)=\frac{1-(1-p_{i,h})^{b_h}}{p_{i,h}},
                \qquad \mathbb E_iW_h=p_{i,h}^{-1}.
\end{equation}
Thus the critical deterministic-cap order for this binary
first-acceptance experiment is
$p_{0,h}^{-1}\asymp h^{-2}e^{\gamma_0J_h}$; vanishing error requires
a diverging multiple of that scale.  This is a scale in preparations,
not the $h^{-8}$ scale in successful marks.
\end{enumerate}
All tests here are for the two specified alternatives; their thresholds
may use the corresponding acceptance probabilities.  The theorem does not
assert a uniform estimator over unknown recording profiles, or an optimal
preparation bound for experiments retaining additional failed-shot data.
\end{theorem}

\begin{proof}
The contact curvatures of the pair are $80/7$ and $40/3$.  Hence
$\gamma_0=\operatorname{arcosh}(71/7)$ and
$\gamma_1=\operatorname{arcosh}(35/3)>\gamma_0$.
The positive fixed-offset prefactor in Theorem~\ref{thm:v4-law} gives
$\pi_{s_i,J}\asymp e^{-\gamma_iJ}$ along this even subsequence.
The bound \eqref{eq:v53-recording-mass} is uniform in small $h$ and large
$J$.  Multiplication proves \eqref{eq:v54-escape-separation}; in
particular both $p_{i,h}$ tend to zero.  These are the original bridge
probabilities and the fixed recording profiles of actual tables, not
freely assigned geometric-distribution parameters.

Independence of the preparations gives the exact factorization, for
$k\ge1$ and measurable $A\subset[-1,1]^2$,
\begin{equation}\label{eq:v54-stopped-factorization}
 \Pr_i(W_h=k,Y_h\in A)
   =(1-p_{i,h})^{k-1}p_{i,h}Q_{s_i,J_h,h}(A).
\end{equation}
Thus the waiting count is geometric on $\{1,2,\ldots\}$ and is
independent of its terminal mark under each alternative.  Summing
\eqref{eq:v54-stopped-factorization} in $k$ shows that discarding the
count gives exactly the finite-flight conditional mark, without any
approximation or selection of a favorable history.
By \eqref{eq:v53-tv} and \eqref{eq:v53-finite-flight},
\begin{equation}\label{eq:v54-mark-distance}
 \operatorname{TV}(E_{0,h},E_{1,h})\le C(h^4+\tau^{J_h})\longrightarrow0.
\end{equation}
Choose $m_h=\lceil(p_{0,h}p_{1,h})^{-1/2}\rceil$ and select alternative
zero when $W_h\le m_h$.  The two errors satisfy
\begin{equation}\label{eq:v54-waiting-errors}
 \Pr_0(W_h>m_h)\le e^{-\rho_h^{-1/2}},\qquad
 \Pr_1(W_h\le m_h)\le \sqrt{\rho_h}+p_{1,h}.
\end{equation}
The first bound uses $(1-p)^m\le e^{-pm}$; the second uses
$1-(1-p)^m\le mp$ and the upper rounding bound for $m_h$.
Both errors vanish, so $\operatorname{TV}(F_{0,h},F_{1,h})\to1$.
The equal-prior error identity proves the first comparison.

Projection discarding $W_h$ proves
$\delta(\mathcal F_h,\mathcal E_h)=0$.  For any kernel $K$, contraction
and the triangle inequality give
\[
 \operatorname{TV}(F_{0,h},F_{1,h})
 \le 2\max_i\operatorname{TV}(KE_{i,h},F_{i,h})
              +\operatorname{TV}(E_{0,h},E_{1,h}).
\]
Equations \eqref{eq:v54-mark-distance}--\eqref{eq:v54-waiting-errors}
therefore give the lower limit $1/2$ for the reverse deficiency.
The kernel ignoring its input and drawing from
$(F_{0,h}+F_{1,h})/2$ has maximum error
$\operatorname{TV}(F_{0,h},F_{1,h})/2\le1/2$ and proves the matching
upper bound.  Kernels may depend on the specified pair and on $h$, but
not on which alternative generated the input.

For the deterministic cap, put
$a_{i,h}=1-(1-p_{i,h})^{b_h}$, the probability of at least one acceptance.
Eventually $p_{1,h}<p_{0,h}$, so $a_{1,h}\le a_{0,h}$.
The two capped laws have respective cemetery masses $1-a_{i,h}$.
Testing the event of any acceptance gives a lower bound on their total
variation, and their common cemetery mass gives an upper bound:
\begin{equation}\label{eq:v54-capped-TV-sandwich}
 a_{0,h}-a_{1,h}
  \le \operatorname{TV}(F_{0,h}^{[b_h]},F_{1,h}^{[b_h]})
  \le a_{0,h}.
\end{equation}
This bound does not presume equality of the conditional accepted-mark
laws.  If $b_hp_{0,h}\to\lambda<\infty$, then
$b_hp_{1,h}=\rho_hb_hp_{0,h}\to0$ and
$b_hp_{0,h}^2\to0$.  Expanding $\log(1-p_{0,h})$ shows
$a_{0,h}\to1-e^{-\lambda}$, while $a_{1,h}\to0$.
The sandwich proves \eqref{eq:v54-capped-critical-risk}.  The stated
test has exact equal-prior error
$\{(1-p_{0,h})^{b_h}+1-(1-p_{1,h})^{b_h}\}/2$ and hence attains that
limit.

If $b_hp_{0,h}\to\infty$, use the earlier test with the observable
threshold $t_h=\min\{b_h,m_h\}$.  It is determined by the capped
record, including its cemetery outcome.  Its errors are bounded by
$e^{-p_{0,h}t_h}$ and $p_{1,h}t_h$.  The first vanishes because
$p_{0,h}t_h=\min\{b_hp_{0,h},p_{0,h}m_h\}\to\infty$; the second is
at most $\sqrt{\rho_h}+p_{1,h}\to0$.  Conversely, if
$b_hp_{0,h}$ does not tend to infinity, a bounded subsequence has a
further subsequence converging to some finite $\lambda$.  On that
subsequence \eqref{eq:v54-capped-critical-risk} gives a strictly
positive limiting error.  This proves the necessity as well as
sufficiency of the cap criterion.

Finally the tail-sum identity gives
$\mathbb E_i\min(W_h,b_h)=\sum_{k=0}^{b_h-1}(1-p_{i,h})^k$,
which is \eqref{eq:v54-capped-charge}; the infinite geometric sum gives
$\mathbb E_iW_h$.  Every failed preparation before stopping is included.
The cap scale follows from \eqref{eq:v54-escape-separation} and the
proved risk criterion, not from turning an expected waiting time into
a deterministic bound.
\end{proof}

The two limits describe distinct coarsenings of one physical acquisition
record.  Conditional normalization suppresses the accepted mass, leaving
an $h^4$ interaction contrast; retaining failures exposes the different
hyperbolic exponents.  At a fixed preparation cap, the cemetery outcome
itself becomes informative.  None of these operations changes the exact
window invariant or the geometric rigidity theorem.  Their statistical
conclusions remain specific to the stated records and fixed pair of
profiles, which need not be known or shared in the unrestricted nuisance
problem.
